\section{AIME I 2012}
        \begin{enumerate}
        	\item (\textit{AIME I 2012}) Find the number of positive integers with three not necessarily distinct digits, $abc$, with $a \neq 0$ and $c \neq 0$ such that both $abc$ and $cba$ are multiples of $4$.


 
 \textbf{Answer: }040



\textbf{Solution 1}

A positive integer is divisible by $4$ if and only if its last two digits are divisible by $4.$ For any value of $b$, there are two possible values for $a$ and $c$, since we find that if $b$ is even, $a$ and $c$ must be either $4$ or $8$, and if $b$ is odd, $a$ and $c$ must be either $2$ or $6$. There are thus $2 \cdot 2 = 4$ ways to choose $a$ and $c$ for each $b,$ and $10$ ways to choose $b$ since $b$ can be any digit. The final answer is then $4 \cdot 10 = \boxed{040}$.

 



\textbf{Solution 2}

A number is divisible by four if its last two digits are divisible by 4. Thus, we require that $10b + a$ and $10b + c$ are both divisible by $4$. If $b$ is odd, then $a$ and $c$ must both be $2 \pmod 4$ meaning that $a$ and $c$ are $2$ or $6$. If $b$ is even, then $a$ and $c$ must be $0 \pmod 4$ meaning that $a$ and $c$ are $4$ or $8$. For each choice of $b$ there are $2$ choices for $a$ and $2$ for $c$ for a total of $10 \cdot 2 \cdot 2 = \boxed{040}$ numbers.

 



\textbf{Solution 3}

For this number to fit the requirements $bc$ and $ba$ must be divisible by 4. So $bc = 00, 04, 08, 12, 16, ... , 92, 96$ and so must $ba$ for each two digits of $bc$. There are two possibilities for $ba$ if $b$ is odd and three possibilities if $b$ is even. So there are $2^{2} \cdot 5+3^{2} \cdot 4 = 65$ possibilities but this overcounts when $a$ or $c = 0$. So when $bc = 00, 20, 40, 60, 80$ and the corresponding $ba$ should be removed, so $65 - 5 \cdot 3 = 50$. But we are still overcounting when $b$ is even because then $a$ can be 0. So the answer is $50 - 10 \cdot 1 = \boxed{040}$ 

 


	\item (\textit{AIME I 2012}) The terms of an arithmetic sequence add to $715$. The first term of the sequence is increased by $1$, the second term is increased by $3$, the third term is increased by $5$, and in general, the $k$th term is increased by the $k$th odd positive integer. The terms of the new sequence add to $836$. Find the sum of the first, last, and middle term of the original sequence.


 
 \textbf{Answer: }195



\textbf{Solution 1}

If the sum of the original sequence is $\sum_{i=1}^{n} a_i$ then the sum of the new sequence can be expressed as $\sum_{i=1}^{n} a_i + (2i - 1) = n^2 + \sum_{i=1}^{n} a_i.$ Therefore, $836 = n^2 + 715 \rightarrow n=11.$ Now the middle term of the original sequence is simply the average of all the terms, or $\frac{715}{11} = 65,$ and the first and last terms average to this middle term, so the desired sum is simply three times the middle term, or $\boxed{195}.$

 
Alternatively, notice that in the original sequence, $11a_1 + 55d = 715$, from which $a_1 + 5d = 65$. Since we are tasked to find $a_1 + a_6 + a_{11} = 3(a_1 + 5d)$, the desired answer is $3 \cdot 65 = \boxed{195}.$

 



\textbf{Solution 2}

After the adding of the odd numbers, the total of the sequence increases by $836 - 715 = 121 = 11^2$. Since the sum of the first $n$ positive odd numbers is $n^2$, there must be $11$ terms in the sequence, so the mean of the sequence is $\dfrac{715}{11} = 65$. Since the first, last, and middle terms are centered around the mean, our final answer is $65 \cdot 3 = \boxed{195}$

 



\textbf{Solution 3}

Proceed as in Solution 2 until it is noted that there are 11 terms in the sequence. Since the sum of the terms of the original arithmetic sequence is 715, we note that $\frac{2a_1 + 10d}{2} \cdot 11 = 715$ or $2a_1 + 10d = 130$ for all sets of first terms and common differences that fit the conditions given in the problem. Assume WLOG that $a_1 = 60$ and $d = 1$. Then the first term of the corresponding arithmetic sequence will be $60$, the sixth (middle) term will be $65$, and the eleventh (largest) term will be $70$. Thus, our final answer is $60 + 65 + 70 = \boxed{195}$.

 



\textbf{Solution 4 (Similar to the above solutions)}

Let the sequence be $a,a+d,a+2d,..., a+(k-1)d$. Since there are $k$ terms, we have an equation $k^2 = 836-715 = 121$. Solving, we get $k$ as $11$. Replacing $k$ in our sequence with $11$, we get $a,a+d,a+2d,..., a+10d$. The sum of this sequence is equal to $715$, or $\frac{(2a+10d)(11)}{2} = 715$, and by simplifying we get $a+5d=65$. We are asked for the first, middle, and end terms, which are $a$,$a+5d$, and $a+10d$ respectively. Their sum is $3a+15$, or $3 \cdot (a+5d)$. Our desired answer is $65 \cdot 3 = \boxed{195}$.

 


	\item (\textit{AIME I 2012}) Nine people sit down for dinner where there are three choices of meals. Three people order the beef meal, three order the chicken meal, and three order the fish meal. The waiter serves the nine meals in random order. Find the number of ways in which the waiter could serve the meal types to the nine people so that exactly one person receives the type of meal ordered by that person.


 
 \textbf{Answer: }216



\textbf{Solution 1}

Call a beef meal $B,$ a chicken meal $C,$ and a fish meal $F.$ Now say the nine people order meals $\text{BBBCCCFFF}$ respectively and say that the person who receives the correct meal is the first person. We will solve for this case and then multiply by $9$ to account for the $9$ different ways in which the person to receive the correct meal could be picked. Note, this implies that the dishes are indistinguishable, though the people aren't. For example, two people who order chicken are separate, though if they receive fish, there is only 1 way to order them.

 The problem we must solve is to distribute meals $\text{BBCCCFFF}$ to orders $\text{BBCCCFFF}$ with no matches. The two people who ordered $B$'s can either both get $C$'s, both get $F$'s, or get one $C$ and one $F.$ We proceed with casework.

 
                \begin{itemize}
                    \item  If the two $B$ people both get $C$'s, then the three $F$ meals left to distribute must all go to the $C$ people. The $F$ people then get $BBC$ in some order, which gives three possibilities. The indistinguishability is easier to see here, as we distribute the $F$ meals to the $C$ people, and there is only 1 way to order this, as all three meals are the same.\item  If the two $B$ people both get $F$'s, the situation is identical to the above and three possibilities arise. \item  If the two $B$ people get $CF$ in some order, then the $C$ people must get $FFB$ and the $F$ people must get $CCB.$ This gives $2 \cdot 3 \cdot 3 = 18$ possibilities. 
                \end{itemize}


                Summing across the cases we see there are $24$ possibilities, so the answer is $9 \cdot 24 = \boxed{216.}$

 



\textbf{Solution 2}

First, choose the person that gets the correct meal. There's obviously $9$ ways to do that. Then, we can casework on who gets what type of meal. WLOG the person who had the correct meal had the chicken meal. The remaining chicken meals can be distributed to either one person who ordered beef and another who ordered fish, or can be distributed to two people who ordered the same type of meal that is not chicken. For the first case, it's apparent that there are $9$ ways of choosing who gets the two chicken meal, and two ways to order the remaining meals. For the second case, there are $6$ ways to choose who gets the two chicken meals and only $1$ way of ordering the remaining meals. Therefore, the answer is $9 \cdot (9 \cdot 2 + 6 \cdot 1) = \boxed {216}$.

 


	\item (\textit{AIME I 2012}) Butch and Sundance need to get out of Dodge. To travel as quickly as possible, each alternates walking and riding their only horse, Sparky, as follows. Butch begins by walking while Sundance rides. When Sundance reaches the first of the hitching posts that are conveniently located at one-mile intervals along their route, he ties Sparky to the post and begins walking. When Butch reaches Sparky, he rides until he passes Sundance, then leaves Sparky at the next hitching post and resumes walking, and they continue in this manner. Sparky, Butch, and Sundance walk at $6,$ $4,$ and $2.5$ miles per hour, respectively. The first time Butch and Sundance meet at a milepost, they are $n$ miles from Dodge, and they have been traveling for $t$ minutes. Find $n + t$.


 
 \textbf{Answer: }279



\textbf{Solution}

When they meet at the milepost, Sparky has been ridden for $n$ miles total. Assume Butch rides Sparky for $a$ miles, and Sundance rides for $n-a$ miles. Thus, we can set up an equation, given that Sparky takes $\frac{1}{6}$ hours per mile, Butch takes $\frac{1}{4}$ hours per mile, and Sundance takes $\frac{2}{5}$ hours per mile:

 
 \[\frac{a}{6} + \frac{n-a}{4} = \frac{n-a}{6} + \frac{2a}{5} \rightarrow  a = \frac{5}{19}n.\]

 The smallest possible integral value of $n$ is $19$, so we plug in $n = 19$ and $a = 5$ and get $t = \frac{13}{3}$ hours, or $260$ minutes. So our answer is $19 + 260 = \boxed{279}$.

 Note that this solution is not rigorous because it is not guaranteed that they will switch properly to form this combination.

 


	\item (\textit{AIME I 2012}) Let $B$ be the set of all binary integers that can be written using exactly $5$ zeros and $8$ ones where leading zeros are allowed. If all possible subtractions are performed in which one element of $B$ is subtracted from another, find the number of times the answer $1$ is obtained.


 
 \textbf{Answer: }330



\textbf{Solution}

When $1$ is subtracted from a binary number, the number of digits will remain constant if and only if the original number ended in $10.$ Therefore, every subtraction involving two numbers from $B$ will necessarily involve exactly one number ending in $10.$ To solve the problem, then, we can simply count the instances of such numbers. With the $10$ in place, the seven remaining $1$'s can be distributed in any of the remaining $11$ spaces, so the answer is ${11 \choose 7} = \boxed{330}$.

 


	\item (\textit{AIME I 2012}) The complex numbers $z$ and $w$ satisfy $z^{13} = w,$ $w^{11} = z,$ and the imaginary part of $z$ is $\sin{\frac{m\pi}{n}}$, for relatively prime positive integers $m$ and $n$ with $m<n.$ Find $n.$


 
 \textbf{Answer: }071



\textbf{Solution}

Substituting the first equation into the second, we find that $(z^{13})^{11} = z$ and thus $z^{143} = z.$ We know that $z \neq 0,$ since we are given the imaginary part of $z$, so we can divide by $z$ to get $z^{142} = 1.$ So, $z$ must be a $142$nd root of unity, and thus, by De Moivre's theorem, the imaginary part of $z$ will be of the form $\sin{\frac{2k\pi}{142}} = \sin{\frac{k\pi}{71}},$ where $k \in \{1, 2, \ldots, 70\}.$ Note that $71$ is prime and $k<71$ by the conditions of the problem, so the denominator in the argument of this value will always be $71.$ Thus, $n = \boxed{071}.$

 



\textbf{Solution 2}

This is similar to solution 1.

 As before, determine $z^{143} = z$, so $z^{142} = 1 \implies \Im(z^{142}) = 0$. So, $\sin(142\theta) = 0$, which can only happen if $142\theta = 2k\pi$ for some $k$. So, $\theta = \frac{k\pi}{71}$. 71 is prime, so $k$ is relatively prime to all $k < 71$ through Euler Totient. Note $\Im(z) = \sin(\theta)$, so the requested imaginary part is $\sin(\frac{k\pi}{71})$, and $n = \boxed{71}$.

 



\textbf{Solution 3}

Giving the same exponent of $w$ gives $z^{143}=w^{11}=z \Longrightarrow z^{143}=z$. We know that $z \neq 0$ from the imaginary part. Thus, $z^{142}=1$, so $z$ is a $142\text{nd}$ root of unity. So, $z=e^{\frac{2\pi i k}{142}}=e^{\frac{\pi i k}{71}}.$ Thus, comparing imaginary parts, we get,
 \[\sin{\dfrac{k\pi}{71}}=\sin{\dfrac{m\pi}{n}} \Longrightarrow \dfrac{k}{71}=\dfrac{m}{n} \Longrightarrow n=\boxed{071}.\]

 -jb2015007

 


	\item (\textit{AIME I 2012}) At each of the sixteen circles in the network below stands a student. A total of $3360$ coins are distributed among the sixteen students. All at once, all students give away all their coins by passing an equal number of coins to each of their neighbors in the network. After the trade, all students have the same number of coins as they started with. Find the number of coins the student standing at the center circle had originally.


 
                    \begin{figure}[H]
                        \centering
                        \includegraphics[width=0.5\textwidth]{Problem/AIME_I_2012_7_Problem_Figure_1.png}
                    \end{figure}
                    

 
 \textbf{Answer: }280



\textbf{Solution 1}

Say the student in the center starts out with $a$ coins, the students neighboring the center student each start with $b$ coins, and all other students start out with $c$ coins. Then the $a$-coin student has five neighbors, all the $b$-coin students have three neighbors, and all the $c$-coin students have four neighbors.

 Now in order for each student's number of coins to remain equal after the trade, the number of coins given by each student must be equal to the number received, and thus

 
 \begin{align*} a &= 5 \cdot \frac{b}{3}\\ b &= \frac{a}{5} + 2 \cdot \frac{c}{4}\\ c &= 2 \cdot \frac{c}{4} + 2 \cdot \frac{b}{3}. \end{align*}


 Solving these equations, we see that $\frac{a}{5} = \frac{b}{3} = \frac{c}{4}.$ Also, the total number of coins is $a + 5b + 10c = 3360,$ so $a + 5 \cdot \frac{3a}{5} + 10 \cdot \frac{4a}{5} = 3360 \rightarrow a = \frac{3360}{12} = \boxed{280}.$

 
                \begin{itemize}
                    \item One way to make this more rigorous (answering the concern in Solution 3) is to let $a$, $b$, and $c$ represent the $total$ number of coins in the rings. Then, you don't care about individual people in the passing, you just know that each ring gets some coins and loses some coins, which must cancel each other out.
                \end{itemize}


                



\textbf{Solution 2}

Since the students give the same number of gifts of coins as they receive and still end up the same number of coins, we can assume that every gift of coins has the same number of coins. Let $x$ be the number of coins in each gift of coins. There $10$ people who give $4$ gifts of coins, $5$ people who give $3$ gifts of coins, and $1$ person who gives $5$ gifts of coins. Thus,

 
 \begin{align*} 10(4x)+5(3x)+5x &= 3360\\ 40x+15x+5x &= 3360\\ 60x &= 3360\\ x &= 56 \end{align*}
Therefore the answer is $5(56) = \boxed{280}.$

 



\textbf{Solution 3}

Mark the number of coins from inside to outside as $a$, $b_1$, $b_2$, $b_3$, $b_4$, $b_5$, $c_1$, $c_2$, $c_3$, $c_4$, $c_5$, $d_1$, $d_2$, $d_3$, $d_4$, $d_5$.Then, we obtain
 \begin{align*} d_1 &= \frac{d_5}{4} + \frac{d_2}{4} + \frac{c_1}{4} + \frac{c_2}{4}\\ d_2 &= \frac{d_1}{4} + \frac{d_3}{4} + \frac{c_2}{4} + \frac{c_3}{4}\\ d_3 &= \frac{d_2}{4} + \frac{d_4}{4} + \frac{c_3}{4} + \frac{c_4}{4}\\ d_4 &= \frac{d_3}{4} + \frac{d_5}{4} + \frac{c_4}{4} + \frac{c_5}{4}\\ d_5 &= \frac{d_4}{4} + \frac{d_1}{4} + \frac{c_5}{4} + \frac{c_1}{4}\\ \end{align*}
Letting $D = d_1 + d_2 + d_3 + d_4 + d_5$, $C = c_1 + c_2 + c_3 + c_4 + c_5$ gets us $D = \frac{D}{4} + \frac{D}{4} + \frac{C}{4} + \frac{C}{4}$ and $D = C$.In the same way, $C = \frac{D}{4} + \frac{D}{4} + \frac{B}{3} + \frac{B}{3}$, $B = \frac{3D}{4}$, $B = \frac{C}{4} + \frac{C}{4} + a$, $a = \frac{D}{4}$.Then, with $a + B + C + D = 3360$, $D = 1120$, $a = \boxed{280}$.

 



\textbf{Solution 4}

Define $x$ as the number of coins the student in the middle has. Since this student connects to $5$ other students, each of those students must have passed $\dfrac15 x$ coins to the center to maintain the same number of coins.

 Each of these students connect to $3$ other students, passing $\dfrac15 x$ coins to each, so they must have $\dfrac35 x$ coins. These students must then recieve $\dfrac35 x$ coins, $\dfrac 15 x$ of which were given to by the center student. Thus, they must also have received $\dfrac25 x$ coins from the outer layer, and since this figure has symmetry, these must be the same.

 Each of the next layer of students must have given $\dfrac15 x$ coins. Since they gave $4$ people coins, they must have started with $\dfrac45 x$ coins. They received $\dfrac25 x$ of them from the inner layer, making the two other connections have given them the same amount. By a similar argument, they recieved $\dfrac15 x$ from each of them.

 By a similar argument, the outermost hexagon of students must have had $\dfrac45 x$ coins each. Summing this all up, we get the number of total coins passed out as a function of $x$. This ends up to be 
 \begin{align*}  \dfrac45 x + \dfrac45 x + \dfrac45 x + \dfrac45 x + \dfrac45 x + \dfrac45 x + \dfrac45 x + \dfrac45 x + \dfrac45 x + \dfrac45 x + \dfrac35 x + \dfrac35 x + \dfrac35 x + \dfrac35 x + \dfrac35 x + x &= 10 \cdot \dfrac45 x  + 5 \cdot \dfrac35 x + x \\  &= 8x + 3x + x \\  &= 12 x.  \end{align*}


 Since this all sums up to $3360$, which is given, we find that 
 \begin{align*} 12x &= 3360 \\ x &= 280 \end{align*}


 We defined $x$ to be the number of coins that the center person has, so the answer is $x$, which is $\boxed{280}$

 


	\item (\textit{AIME I 2012}) Cube $ABCDEFGH,$ labeled as shown below, has edge length $1$ and is cut by a plane passing through vertex $D$ and the midpoints $M$ and $N$ of $\overline{AB}$ and $\overline{CG}$ respectively. The plane divides the cube into two solids. The volume of the larger of the two solids can be written in the form $\frac{p}{q},$ where $p$ and $q$ are relatively prime positive integers. Find $p+q.$


 
                    \begin{figure}[H]
                        \centering
                        \includegraphics[width=0.5\textwidth]{Problem/AIME_I_2012_8_Problem_Figure_1.png}
                    \end{figure}
                    

 
 \textbf{Answer: }089



\textbf{Solution: think outside the box (pun intended)}

Define a coordinate system with $D$ at the origin and $C,$ $A,$ and $H$ on the $x$, $y$, and $z$ axes respectively. Then $D=(0,0,0),$ $M=(.5,1,0),$ and $N=(1,0,.5).$ It follows that the plane going through $D,$ $M,$ and $N$ has equation $2x-y-4z=0.$ Let $Q = (1,1,.25)$ be the intersection of this plane and edge $BF$ and let $P = (1,2,0).$ Now since $2(1) - 1(2) - 4(0) = 0,$ $P$ is on the plane. Also, $P$ lies on the extensions of segments $DM,$ $NQ,$ and $CB$ so the solid $DPCN = DMBCQN + MPBQ$ is a right triangular pyramid. Note also that pyramid $MPBQ$ is similar to $DPCN$ with scale factor $.5$ and thus the volume of solid $DMBCQN,$ which is one of the solids bounded by the cube and the plane, is $[DPCN] - [MPBQ] = [DPCN] - \left(\frac{1}{2}\right)^3[DPCN] = \frac{7}{8}[DPCN].$ But the volume of $DPCN$ is simply the volume of a pyramid with base $1$ and height $.5$ which is $\frac{1}{3} \cdot 1 \cdot .5 = \frac{1}{6}.$ So $[DMBCQN] = \frac{7}{8} \cdot \frac{1}{6} = \frac{7}{48}.$ Note, however, that this volume is less than $.5$ and thus this solid is the smaller of the two solids. The desired volume is then $[ABCDEFGH] - [DMBCQN] = 1 - \frac{7}{48} = \frac{41}{48} \rightarrow p+q = \boxed{089.}$

 
                \begin{itemize}
                    \item Another way to finish after using coordinates: Take the volume of DMBCNQ as the sum of the volumes of DMBQ and DBCNQ. Then, we have the sum of the volumes of a tetrahedron and a pyramid with a trapezoidal base. Their volumes, respectively, are $\frac{1}{48}$ and $\frac{1}{8}$, giving the same answer as above. $\blacksquare$ ~mathtiger6
                \end{itemize}


                



\textbf{Solution 2}

Define a coordinate system with $D = (0,0,0)$, $M = (1, \frac{1}{2}, 0)$, $N = (0,1,\frac{1}{2})$. The plane formed by $D$, $M$, and $N$ is $z = \frac{y}{2} - \frac{x}{4}$. It intersects the base of the unit cube at $y = \frac{x}{2}$. The z-coordinate of the plane never exceeds the height of the unit cube for $0 \leq x \leq 1, 0 \leq y \leq 1$. Therefore, the volume of one of the two regions formed by the plane is
 \[\int_0^1 \int_{\frac{x}{2}}^1 \int_0^{\frac{y}{2}-\frac{x}{4}}dz\,dy\,dx = \frac{7}{48}\]Since $\frac{7}{48} < \frac{1}{2}$, our answer is $1-\frac{7}{48} = \frac{41}{48} \rightarrow p+q = \boxed{089}$.

 



\textbf{Solution 2 but slightly better}

Same coordinate system, but we will use domains instead of using triple integrals. By the way, the method to obtain the equation of the plane is the cross product ($\vec{[1,.5,0]}\times\vec{[0,1,.5]}=\vec{[.25,-.5,1]}$). We can multiply this vector by $-4$ to make things look cleaner and get $\vec{[-1,2,-4]}$. We then get the desired plane, $-x+2y-4z=0$, or $z=\frac{2y-x}{4}$. We use a double integral with a Type I domain. Observing the diagram, the domain is where $0 \leq x \leq 1, .5x \leq y \leq 1$. The integral is then
 \[.25\int_0^1\int_{.5x}^1 2y - x\]which becomes
 \[.25\int_0^1 (y^2-xy)|_{.5x}^1\]which becomes
 \[.25\int_0^1 1-x+.25x^2\]which then becomes
 \[.25(1-\frac{1}{2}+\frac{1}{12})\]and finally
 \[\frac{7}{48}\]So our answer is $1^3-\frac{7}{48} = \frac{41}{48}, 41+48 = \boxed{089}$

 \textbf{Alternative Solution (using calculus) : think inside the box} Let $Q$ be the intersection of the plane with edge $FB,$ then $\triangle MQB$ is similar to $\triangle DNC$ and the volume $[DNCMQB]$ is a sum of areas of cross sections of similar triangles running parallel to face $ABFE.$ Let $x$ be the distance from face $ABFE,$ let $h$ be the height parallel to $AB$ of the cross-sectional triangle at that distance, and $a$ be the area of the cross-sectional triangle at that distance. $A(x)=\frac{h^2}{4},$ and $h=\frac{x+1}{2},$ then $A=\frac{(x+1)^2}{16}$, and the volume $[DNCMQB]$ is $\int^1_0{A(x)}\,\mathrm{d}x=\int^1_0{\frac{(x+1)^2}{16}}\,\mathrm{d}x=\frac{7}{48}.$ Thus the volume of the larger solid is $1-\frac{7}{48}=\frac{41}{48} \rightarrow p+q = \boxed{089}$

 \textbf{Alternative Solution : think inside the box} The volume of a frustum is $\frac{h_2b_2 -h_1b_1}3$ where $b_i$ is the area of the base and $h_i$ is the height from the chopped off apex to the base.

 We can easily see that from symmetry, the area of the smaller front base is $\frac{1}{16}$ and the area of the larger back base is $\frac{1}4$

 Now to find the height of the apex.

 Extend the $DM$ and (call the intersection of the plane with $FB$ G) $NG$ to meet at $x$. Now from similar triangles $XMG$ and $XDN$ we can easily find the total height of the triangle $XDN$ to be $2$

 Now straight from our formula, the volume is $\frac{7}{48}$ Thus the answer is:

 $1-\text{Volume} \Longrightarrow \boxed{089}$

 \textbf{Alternative Solution : think inside the box with pyramids} We will solve for the area of the smaller region, and then subtract it from 1.

 Let $X$ be the point where plane $DMN$ intersects $FB$. Then $DMBCNX$ can be split into triangular pyramid $DMBX$ and quadrilateral pyramid $BCNXD$.

 Pyramid $DMBX$ has base $DMB$ with area $\frac{1}2 \cdot 1 \div 2 = \frac{1}4$. The height is $BX = \frac{1}4$, so the volume of $DMBX$ is $\frac{1}4 \cdot \frac{1}4 \div 3 = \frac{1}{48}$.

 Similarly, pyramid $BCNXD$ has base $BCNX$ with area $(\frac{1}4 + \frac{1}2) \cdot 1 = \frac{3}8$. The height is $CD = 1$, so the volume of $BCNXD$ is $\frac{3}8 \cdot 1 \div 3 = \frac{1}8$.

 Adding up the volumes of $DMBX$ and $BCNXD$, we find that the volume of $DMBCNX$ is $\frac{1}{48} + \frac{6}{48} = \frac{7}{48}$. Therefore the volume of the larger solid is $1 - \frac{7}{48} = \frac{41}{48} \rightarrow p+q = \boxed{089}$

 This is basically mathtiger6's solution, but you don't need coordinates or thinking outside the box.

 



\textbf{Solution outside visual}


                    \begin{figure}[H]
                        \centering
                        \includegraphics[width=0.5\textwidth]{Solution/AIME_I_2012_8_Solution_Figure_1.png}
                    \end{figure}
                    

 



\textbf{Solution 7}

We use a coordinate system with $C = (0,0,0)$, $D = (1,0,0)$, $M = (\frac{1}{2},1,0)$, $N = (0,0,\frac{1}{2})$. Then the plane going through $D$, $M$, and $N$ has equation $z = \frac{1}{2} - \frac{1}{2}x - \frac{1}{4}y$. We set up a double integral in this coordinate system. Consider the region to be integrated over in the $xy$-plane. From $x=0$ to $x=\frac{1}{2}$, the upper bound of the region is $y = 1$. From $x= \frac{1}{2}$ to $x=1$, the upper bound of the region is $y = 2 - 2x$. In both cases, the lower bound of the region is $y = 0$. Thus, we have the double integral $\int_{0}^{\frac{1}{2}} \int_{0}^{1} \frac{1}{2} - \frac{1}{2}x - \frac{1}{4}y dydx + \int_{\frac{1}{2}}^{1} \int_{0}^{2-2x} \frac{1}{2} - \frac{1}{2}x - \frac{1}{4}y dydx$. We find that this sum evaluates to $\frac{7}{48}$. However, this is the volume of the smaller region, so the larger region's volume is that of the cube minus that of the smaller region. Since the cube has side length $1$, its volume is $1$, so the volume of the larger region is $1 - \frac{7}{48} = \frac{41}{48}$. Thus, our answer is $41 + 48 = \boxed{089}$.

 


	\item (\textit{AIME I 2012}) Let $x,$ $y,$ and $z$ be positive real numbers that satisfy

 \[2\log_{x}(2y) = 2\log_{2x}(4z) = \log_{2x^4}(8yz) \ne 0.\]
The value of $xy^5z$ can be expressed in the form $\frac{1}{2^{p/q}},$ where $p$ and $q$ are relatively prime positive integers. Find $p+q.$


 
 \textbf{Answer: }049



\textbf{Solution 1}

Since there are only two dependent equations given and three unknowns, the three expressions given can equate to any common value (that isn't 0, of course), so to simplify the problem let us assume without loss of generality that
 \[2\log_{x}(2y) = 2\log_{2x}(4z) = \log_{2x^4}(8yz) = 2.\]Then
 \begin{align*} 2\log_{x}(2y) = 2 &\implies x=2y\\ 2\log_{2x}(4z) = 2 &\implies 2x=4z\\ \log_{2x^4}(8yz) = 2 &\implies 4x^8 = 8yz \end{align*}
Solving these equations, we quickly see that $4x^8 = (2y)(4z) = x(2x) \rightarrow x=2^{-1/6}$ and then $y=z=2^{-1/6 - 1} = 2^{-7/6}.$Finally, our desired value is $2^{-1/6} \cdot (2^{-7/6})^5 \cdot 2^{-7/6} = 2^{-43/6}$ and thus $p+q = 43 + 6 = \boxed{049.}$

 



\textbf{Solution 2}

Notice that $2y\cdot 4z=8yz$, $2\log_2(2y)=\log_2\left(4y^2\right)$ and $2\log_2(4z)=\log_2\left(16z^2\right)$.

 From this, we see that $8yz$ is the geometric mean of $4y^2$ and $16z^2$. So, for constant $C\ne 0$:
 \[\frac{\log 4y^2}{\log x}=\frac{\log 8yz}{\log 2x^4}=\frac{\log 16z^2}{\log 2x}=C\]Since $\log 4y^2,\log 8yz,\log 16z^2$ are in an geometric progression (ratio $\frac{2z}{y}$), so are $\log x,\log 2x^4,\log 2x$.

 Therefore, $2x^4$ is the geometric mean of $x$ and $2x$
 \[2x^4=\sqrt{x\cdot 2x}\implies 4x^8=2x^2\implies 2x^6=1\implies x=2^{-1/6}\]We can plug $x$ in to any of the two equal fractions aforementioned. So, without loss of generality:
 \[\frac{\log 4y^2}{\log x}=\frac{\log 16z^2}{\log 2x}\implies \log\left(4y^2\right)\log\left(2x\right)=\log\left(16z^2\right)\log\left(x\right)\]
 \[\implies \log\left(4y^2\right)\cdot \frac{5}{6}\log 2=\log\left(16z^2\right)\cdot \frac{-1}{6}\log 2\]
 \[\implies 5\log\left(4y^2\right)=-\log\left(16z^2\right)\implies 5\log\left(4y^2\right)+\log\left(16z^2\right)=0\]
 \[\implies \left(4y^2\right)^5\cdot 16z^2=1\implies 16384y^{10}z^2=1\implies y^{10}z^2=\frac{1}{16384}\implies y^5z=\frac{1}{128}\]

 Thus $xy^5z=2^{-\frac{1}{6}-7}=2^{-\frac{43}{6}}$ and $43+6=\boxed{049}$.

 



\textbf{Solution 3}

Since $x, y, z$ are positive real numbers, $x, y$, and $z$ can be expressed as powers of two.  We shall thus let $x = 2^X$, $y = 2^Y$, and $z = 2^Z$.  Let $a = \log_{2^X}(2^{Y+1})$.  From this we get the system of equations: 
 \[\]$(1)$
 \[a = \log_{2^X}(2^{Y+1}) \Rightarrow aX = Y + 1\]$(2)$
 \[a = \log_{2^{X + 1}}(2^{Z + 2}) \Rightarrow aX + a = Z + 2\]$(3)$
 \[2a = \log_{2^{4X + 1}}(2^{Y + Z + 3}) \Rightarrow 8aX + 2a = Y + Z + 3\]

 Plugging equation $(1)$ into equation $(2)$ yields $Y + a = Z + 1$.  Plugging equation $(1)$ into equation $(3)$ and simplifying yields $7Y + 2a + 6 = Z + 1$, and substituting $Y + a$ for $Z + 1$ and simplifying yields $Y + 1 = \frac{-a}{6}$.  But $Y + 1 = aX$, so $aX =  \frac{-a}{6}$, so $X = \frac{-1}{6}$.

 Knowing this, we may substitute $\frac{-1}{6}$ for $X$ in equations $(1)$ and $(2)$, yielding $\frac{-a}{6} = Y + 1$ and $\frac{5a}{6} = Z + 2$.  Thus, we have that $-5(Y + 1) = Z + 2 \rightarrow 5Y + Z = -7$.  We are looking for $xy^5z = 2^{X+ 5Y + Z}$. $X = \frac{-1}{6}$ and $5Y + Z = -7$, so $xy^5z = 2^{-43/6} = \frac{1}{2^{43/6}}$.  The answer is $43+6=\boxed{049}$.

 



\textbf{Solution 4 (Rigorous and easy)}

By the \href{https://artofproblemsolving.com/wiki/index.php?title=Mediant\_theorem}{Mediant theorem}, we know that
 \[\frac{\log (4y^2)}{\log (x)} = \frac{\log (16z^2)}{\log (2x)} = \frac{2\log (8yz)}{\log (2x^2)}.\] 

 Substituting into the original equation yields us $\frac{2\log (8yz)}{\log (2x^2)} = \frac{\log (8yz)}{\log (2x^4)} \Rightarrow 2\log (2x^4) = \log (2x^2) \Rightarrow x=2^{-1/6}.$For some constant $C\not= 0,$ Let $2\log_{x}(2y) = 2\log_{2x}(4z) = \log_{2x^4}(8yz) = C.$ Then, we obtain the system of equations
 \[y=2^{-13C/12}\]
 \[z=2^{-19C/12}\]
 \[8yz=2^{C/3}.\]

 Multiplying the first two equations and dividing by the third, we find $C=1.$ Thus, 
 \[xy^5z=2^{-1/6} \cdot 2^{-65/12} \cdot 2^{-19/12}=2^{-43/6} \Rightarrow p+q=\boxed{049}.\]

 



\textbf{Solution 5 (No Fancy Tricks)}

First, we can move the $2$'s inside the logarithms by squaring the stuff inside: $\log_{x}(4y^2) = \log_{2x}(16z^2) = \log_{2x^4}(8yz)$.

 Setting these 3 expressions to equal $a$, we can get rid of the logarithms by using $\log_{a}b = c \implies a^c = b$. The resulting equations are:

 \begin{align}x^a &= 4y^2 & (1) \\2^a x^a &= 16z^2 & (2) \\2^a x^{4a} &= 8yz & (3) \\\end{align}

 Multiplying equations $(1)$ and $(2)$ yields 
 \[2^a x^{2a} = 64y^2z^2\]

 Squaring equation $(3)$ yields 
 \[2^{2a} x^{8a} = 64y^2z^2\]

 Equating yields

 \begin{align}2^a x^{2a} &= 2^{2a} x^{8a} \\2^{-a} &= x^{6a} \\x &= 2^{-\frac{1}{6}} \\\end{align}

 
Dividing equation $(3)$ by equation $(2)$ yields 
 \[x^{3a} = 2^{-1}yz^{-1}\]

 Cubing equation $(1)$ yields 
 \[x^{3a} = 2^6y^6\]

 Equating yields

 \begin{align}2^{-1}yz^{-1} &= 2^6y^6 \\z^{-1} &= 2^7y^5\\z &= 2^{-7}y^{-5}\\\end{align}

 
Thus, $xy^5z = 2^{-\frac{1}{6}} y^5 2^{-7}y^{-5} = 2^{-\frac{43}{6}} = \dfrac{1}{2^{\frac{43}{6}}} \implies p+q = \boxed{49}$.

 


	\item (\textit{AIME I 2012}) Let $\mathcal{S}$ be the set of all perfect squares whose rightmost three digits in base $10$ are $256$. Let $\mathcal{T}$ be the set of all numbers of the form $\frac{x-256}{1000}$, where $x$ is in $\mathcal{S}$. In other words, $\mathcal{T}$ is the set of numbers that result when the last three digits of each number in $\mathcal{S}$ are truncated. Find the remainder when the tenth smallest element of $\mathcal{T}$ is divided by $1000$.


 
 \textbf{Answer: }170



\textbf{Solution 1}

It is apparent that for a perfect square $s^2$ to satisfy the constraints, we must have $s^2 - 256 = 1000n$ or $(s+16)(s-16) = 1000n.$ Now in order for $(s+16)(s-16)$ to be a multiple of $1000,$ at least one of $s+16$ and $s-16$ must be a multiple of $5,$ and since $s+16 \not\equiv s-16 \pmod{5},$ one term must have all the factors of $5$ and thus must be a multiple of $125.$ Furthermore, each of $s+16$ and $s-16$ must have at least two factors of $2,$ since otherwise $(s+16)(s-16)$ could not possibly be divisible by $8.$ So therefore the conditions are satisfied if either $s+16$ or $s-16$ is divisible by $500,$ or equivalently if $s = 500n \pm 16.$ Counting up from $n=0$ to $n=5,$ we see that the tenth value of $s$ is $500 \cdot 5 - 16 = 2484$ and thus the corresponding element in $\mathcal{T}$ is $\frac{2484^2 - 256}{1000} = 6170 \rightarrow \boxed{170.}$

 



\textbf{Solution 2}

Notice that is $16^2=256$, $1016^2$ ends in $256$. In general, if $x^2$ ends in $256$, $(x+1000)^2=x^2+2000x+1000000$ ends in 256 because $1000|2000x$ and $1000|2000000$. It is clear that we want all numbers whose squares end in $256$ that are less than $1000$.

 Firstly, we know the number has to end in a $4$ or a $6$. Let's look at the ones ending in $6$.

 Assume that the second digit of the three digit number is $0$. We find that the last $3$ digits of $\overline{a06}^2$ is in the form $12a \cdot 100 + 3 \cdot 10 + 6$. However, the last two digits need to be a $56$. Thus, similarly trying all numbers from $0$ to $10$, we see that only 1 for the middle digit works. Carrying out the multiplication, we see that the last $3$ digits of $\overline{a16}^2$ is in the form $(12a + 2) \cdot 100 + 5 \cdot 10 + 6$. We want $(12a + 2)\pmod{10}$ to be equal to $2$. Thus, we see that a is $0$ or $5$. Thus, the numbers that work in this case are $016$ and $516$. 

 Next, let's look at the ones ending in $4$. Carrying out a similar technique as above, we see that the last $3$ digits of $\overline{a84}^2$ is in the form $((8a+10) \cdot 100+ 5 \cdot 10 + 6$. We want $(8a + 10)\pmod{10}$ to be equal to $2$. We see that only $4$ and $9$ work. Thus, we see that only $484$ and $984$ work.

 We order these numbers to get  $16$, $484$, $516$, $984$... We want the $10th$ number in order which is $2484^2 = 6\boxed{170}256$.

 



\textbf{Solution 3}

The condition implies $x^2\equiv 256 \pmod{1000}$. Rearranging and factoring, 
 \[(x-16)(x+16)\equiv 0\pmod {1000}.\]This can be expressed with the system of congruences
 \[\begin{cases}  (x-16)(x+16)\equiv 0\pmod{125} \\ (x-16)(x+16)\equiv 0\pmod{8} \end{cases}\]Observe that $x\equiv {109} \pmod {125}$ or $x\equiv{16}\pmod {125}$. Similarly, it can be seen that $x\equiv{0}\pmod{8}$ or $x\equiv{4}\pmod{8}$. By CRT, there are four solutions to this modulo $1000$ (one for each case e.g. $x\equiv{109}$ and $x\equiv{0}$ or $x\equiv{125}$ and $x\equiv{4}$. These solutions are (working modulo $1000$)
 \[\begin{cases}  x=16 \\ x=484 \\ x= 516 \\ x=984 \end{cases}\] The tenth solution is $x=2484,$ which gives an answer of $\boxed{170}$.

 



\textbf{Solution 4}

An element in S can be represented by $y^2 = 1000a + 256$, where $y^2$ is the element in S. Since the right hand side must be even, we let $y = 2y_1$ and substitute to get $y_1^2 = 250a + 64$. However, the right hand side is still even, so we make the substitution $y_1 = 2y_2$ to get $y_2^2 = 125a/2 + 16$. Because both sides must be an integer, we know that $a = 2a_1$ for some integer $a_1$. Our equation then becomes $y_2^2 = 125a_1 + 16$, and we can simplify no further.

 Rearranging terms, we get $y_2^2 - 16 = 125a_1$, whence difference of squares gives $(y_2 + 4)(y_2 - 4) = 125a_1$. Note that this equation tells us that one of $y_2 + 4$ and $y_2 - 4$ contains a nonnegative multiple of $125$. Hence, listing out the smallest possible values of $y_2$, we have $y_2 = 4, 121, 129, \cdots, 621$. The tenth term is $621$, whence $y = 4y_2 = 2484$. The desired result can then be calculated to be $\boxed{170}$. - Spacesam

 



\textbf{Solution 5 (Similar to Solution 4)}

From the conditions, we can let every element in $\mathcal{S}$ be written as $y^2=1000x+256$, where $x$ and $y$ are integers. Since there are no restrictions on $y$, we let $y_1$ be equal to $y+16$ ($y-16$ works as well). Then the $256$ cancels out and we're left with 
 \[y_1^2+32y_1=1000x\] which can be factored as 
 \[y_1(y_1+32)=1000x\] Since the RHS is even, $y_1$ must be even, so we let $y_1=2y_2$, to get 
 \[y_2(y_2+16)=250x\] Again, because the RHS is even, the LHS must be even too, so substituting $y_3=\frac{1}{2}y_2$ we have 
 \[y_3(y_3+8)=125\cdot\frac{1}{2}x\] Since the LHS is an integer, the RHS must thus be an integer, so substituting $x=2x_1$ we get 
 \[y_3(y_3+8)=125x_1\] Then we can do casework on the values of $y_3$, as only one of $y_3$ and $y_3+8$ can be multiples of $125$

 \textbf{Case 1}: $125|y_3$

 Since we're trying to find the values of $x_1$, we can let $y_4=\frac{1}{125}y_3$, to get 
 \[x_1=y_4(125y_4+8)\] or 
 \[x=2y_4(125y_4+8)\]

 \textbf{Case 2}: $125|y_3+8$

 Similar to Case 1, only the equation is 
 \[x=2y_4(125y_4-8)\]

 In whole, the values of $x$ (i.e. the elements in $\mathcal{T}$) are of the form 
 \[x=2k(125k\pm8)\] where $k$ is any integer. It can easily be seen that if $k<0$, then $x$ is negative, thus $k\geq0$. Also, note that when $k=0$, there is only one value, because one of the factors is $0$ ($k$). Thus the $10^{th}$ smallest number in the set $\mathcal{T}$ is when the $\pm$ sign takes the minus side and $k=5$, giving $6170$, so the answer is $\boxed{170}$

 


	\item (\textit{AIME I 2012}) A frog begins at $P_0 = (0,0)$ and makes a sequence of jumps according to the following rule: from $P_n = (x_n, y_n),$ the frog jumps to $P_{n+1},$ which may be any of the points $(x_n + 7, y_n + 2),$ $(x_n + 2, y_n + 7),$ $(x_n - 5, y_n - 10),$ or $(x_n - 10, y_n - 5).$ There are $M$ points $(x, y)$ with $|x| + |y| \le 100$ that can be reached by a sequence of such jumps. Find the remainder when $M$ is divided by $1000.$


 
 \textbf{Answer: }373



\textbf{Solution}

First of all, it is easy to see by induction that for any $P(x,y)$ in the frog's jump sequence, $x+y$ will be a multiple of $3$ and $x-y$ will be a multiple of $5.$ The base case $(x,y) = (0,0)$ obviously satisfies the constraints and if $x+y = 3n$ and $x-y = 5m,$ any of the four transformations will sustain this property:

 
 \begin{align*} (x+7)+(y+2) = x+y+9 \rightarrow 3(n+3) &\text{ and } (x+7)-(y+2) = x-y+5 \rightarrow 5(m+1)\\ (x+2)+(y+7) = x+y+9 \rightarrow 3(n+3) &\text{ and } (x+2)-(y+7) = x-y-5 \rightarrow 5(m-1)\\ (x-5)+(y-10) = x+y-15 \rightarrow 3(n-5) &\text{ and } (x-5)-(y-10) = x-y+5 \rightarrow 5(m+1)\\ (x-10)+(y-5) = x+y-15 \rightarrow 3(n-5) &\text{ and } (x-10)-(y-5) = x-y-5 \rightarrow 5(m-1).\\ \end{align*}
So we know that any point the frog can reach will satisfy $x+y = 3n$ and $x-y = 5m.$ 

 $\textbf{Lemma:}$ Any point $(x,y)$ such that there exists 2 integers $m$ and $n$ that satisfy $x+y = 3n$ and $x-y = 5m$ is reachable. 

 
$\textbf{Proof:}$ Denote the total amounts of each specific transformation in the frog's jump sequence to be $a,$ $b,$ $c,$ and $d$ respectively. Then

 $x=7a+2b-5c-10d$,

 $y=2a+7b-10c-5d$,

 $x+y = 9(a+b)-15(c+d) = 3n$, and

 $x-y = 5(a-b)+5(c-d) = 5m$

 together must have integral solutions. But

 $3(a+b)-5(c+d) = n$ implies

 $(c+d) \equiv n \mod 3$ and thus

 $(a+b) = \lfloor{n/3}\rfloor + 2(c+d).$

 Similarly, $(a-b)+(c-d) = m$ implies that $(a-b)$ and $(c-d)$ have the same parity. Now in order for an integral solution to exist, there must always be a way to ensure that the pairs $(a+b)$ and $(a-b)$ and $(c+d)$ and $(c-d)$ have identical parities. The parity of $(a+b)$ is completely dependent on $n,$ so the parities of $(a-b)$ and $(c-d)$ must be chosen to match this value. But the parity of $(c+d)$ can then be adjusted by adding or subtracting $3$ until it is identical to the parity of $(c-d)$ as chosen before, so we conclude that it is always possible to find an integer solution for $(a,b,c,d)$ and thus any point that satisfies $x+y = 3n$ and $x-y = 5m$ can be reached by the frog.

 To count the number of such points in the region $|x| + |y| \le 100,$ we first note that any such point will lie on the intersection of one line of the form $y=x-5m$ and another line of the form $y=-x+3n.$ The intersection of two such lines will yield the point $\left(\frac{3n+5m}{2},\frac{3n-5m}{2}\right),$ which will be integral if and only if $m$ and $n$ have the same parity. Now since $|x| + |y| = |x \pm y|,$ we find that

 
 \begin{align*} |x + y| = |3n| \le 100 &\rightarrow -33 \le n \le 33\\ |x - y| = |5m| \le 100 &\rightarrow -20 \le m \le 20. \end{align*}


 So there are $34$ possible odd values and $33$ possible even values for $n,$ and $20$ possible odd values and $21$ possible even values for $m.$ Every pair of lines described above will yield a valid accessible point for all pairs of $m$ and $n$ with the same parity, and the number of points $M$ is thus $34 \cdot 20 + 33 \cdot 21 = 1373 \rightarrow \boxed{373}$.

 


	\item (\textit{AIME I 2012}) Let $\triangle ABC$ be a right triangle with right angle at $C.$ Let $D$ and $E$ be points on $\overline{AB}$ with $D$ between $A$ and $E$ such that $\overline{CD}$ and $\overline{CE}$ trisect $\angle C.$ If $\frac{DE}{BE} = \frac{8}{15},$ then $\tan B$ can be written as $\frac{m \sqrt{p}}{n},$ where $m$ and $n$ are relatively prime positive integers, and $p$ is a positive integer not divisible by the square of any prime. Find $m+n+p.$


 
 \textbf{Answer: }018



\textbf{Solution 1}

We have $\angle BCE = \angle ECD = \angle DCA = \frac 13 \cdot 90^\circ = 30^\circ$. Drop the altitude from $D$ to $CB$ and call the foot $F$. 

 
                    \begin{figure}[H]
                        \centering
                        \includegraphics[width=0.5\textwidth]{Solution/AIME_I_2012_12_Solution_Figure_1.png}
                    \end{figure}
                    

 Let $CD = 8a$. Using angle bisector theorem on $\triangle CDB$, we get $CB = 15a$. Now $CDF$ is a $30$-$60$-$90$ triangle, so $CF = 4a$, $FD = 4a\sqrt{3}$, and $FB = 11a$.  Finally, $\tan{B} = \frac{DF}{FB}=\frac{4\sqrt{3}a}{11a} = \frac{4\sqrt{3}}{11}$. Our final answer is $4 + 3 + 11 = \boxed{018}$.

 



\textbf{Solution 2}

Without loss of generality, set $CB = 1$. Then, by the Angle Bisector Theorem on triangle $DCB$, we have $CD = \frac{8}{15}$. We apply the Law of Cosines to triangle $DCB$ to get $1 + \frac{64}{225} - \frac{8}{15} = BD^{2}$, which we can simplify to get $BD = \frac{13}{15}$. 

 Now, we have $\cos \angle B = \frac{1 + \frac{169}{225} - \frac{64}{225}}{\frac{26}{15}}$ by another application of the Law of Cosines to triangle $DCB$, so $\cos \angle B = \frac{11}{13}$. In addition, $\sin \angle B = \sqrt{1 - \frac{121}{169}} = \frac{4\sqrt{3}}{13}$, so $\tan \angle B = \frac{4\sqrt{3}}{11}$. 

 Our final answer is $4+3+11 = \boxed{018}$.

 



\textbf{Solution 3}

(This solution does not use the Angle Bisector Theorem or the Law of Cosines, but it uses the Law of Sines and more trig)

 Find values for all angles in terms of $\angle B$. $\angle CEB = 150-B$, $\angle CED = 30+B$, $\angle CDE = 120-B$, $\angle CDA = 60+B$, and $\angle A = 90-B$.

 Use the law of sines on $\triangle CED$ and $\triangle CEB$:

 In $\triangle CED$, $\frac{8}{\sin 30} = \frac{CE}{\sin (120-B)}$. This simplifies to $16 = \frac{CE}{\sin (120-B)}$.

 In $\triangle CEB$, $\frac{15}{\sin 30} = \frac{CE}{\sin B}$. This simplifies to $30 = \frac{CE}{\sin B}$.

 Solve for $CE$ and equate them so that you get $16\sin (120-B) = 30\sin B$. 

 From this, $\frac{8}{15} = \frac{\sin B}{\sin (120-B)}$. 

 Use a trig identity on the denominator on the right to obtain:  $\frac{8}{15} = \frac{\sin B}{\sin 120 \cos B - \cos 120 \sin B}$

 This simplifies to $\frac{8}{15} = \frac{\sin B}{\frac{\sqrt{3}\cos B}{2} + \frac{\sin B}{2}} = \frac{\sin B}{\frac{\sqrt{3} \cos B + \sin B}{2}} = \frac{2\sin B}{\sqrt{3}\cos B + \sin B}$

 This gives $8\sqrt{3}\cos B+8\sin B=30\sin B$Dividing by $\cos B$, we have ${8\sqrt{3}}= 22\tan B$

 $\tan{B} = \frac{8\sqrt{3}}{22} = \frac{4\sqrt{3}}{11}$. Our final answer is $4 + 3 + 11 = \boxed{018}$.

 



\textbf{Solution 4}

(This solution avoids advanced trigonometry)

 Let $X$ be the foot of the perpendicular from $D$ to $\overline{BC}$, and let $Y$ be the foot of the perpendicular from $E$ to $\overline{BC}$.

 Now let $EY=x$. Clearly, triangles $EYB$ and $DXB$ are similar with $\frac{BE}{BD}=\frac{15}{15+8}=\frac{15}{23}=\frac{EY}{DX}$, so $DX=\frac{23}{15}x$.

 Since triangles $CDX$ and $CEY$ are 30-60-90 right triangles, we can easily find other lengths in terms of $x$. For example, we see that $CY=x\sqrt{3}$ and $CX=\frac{\frac{23}{15}x}{\sqrt{3}}=\frac{23\sqrt{3}}{45}x$. Therefore $XY=CY-CX=x\sqrt{3}-\frac{23\sqrt{3}}{45}x=\frac{22\sqrt{3}}{45}x$.

 Again using the fact that triangles $EYB$ and $DXB$ are similar, we see that $\frac{BX}{BY}=\frac{XY+BY}{BY}=\frac{XY}{BY}+1=\frac{23}{15}$, so $BY=\frac{15}{8}XY=\frac{15}{8}*\frac{22\sqrt{3}}{45}=\frac{11\sqrt{3}}{2}$.

 Thus $\tan \angle B = \frac{x}{\frac{11\sqrt{3}}{12}x}=\frac{4\sqrt{3}}{11}$, and our answer is $4+3+11=\boxed{018}$.

 



\textbf{Solution 5}

(Another solution without trigonometry)

 Extend $CD$ to point $F$ such that $\overline{AF} \parallel \overline{CB}$. It is then clear that $\triangle AFD$ is similar to $\triangle BCD$.

 Let $AC=p$, $BC=q$. Then $\tan \angle B = p/q$.

 With the Angle Bisector Theorem, we get that $CD=\frac{8}{15}q$. From 30-60-90 $\triangle CAF$, we get that $AF=\frac{1}{\sqrt{3}}p$ and $FD=FC-CD=\frac{2}{\sqrt{3}}p-\frac{8}{15}q$.

 From $\triangle AFD \sim \triangle BCD$, we have that $\frac{FD}{CD}=\frac{FA}{CB}=\frac{\frac{2}{\sqrt{3}}p-\frac{8}{15}q}{\frac{8}{15}q}=\frac{\frac{1}{\sqrt{3}}p}{q}$. Simplifying yields $\left(\frac{p}{q}\right)\left(\frac{2\sqrt{3}}{3}*\frac{15}{8}-\frac{\sqrt{3}}{3}\right)=1$, and $\tan \angle B=\frac{p}{q}=\frac{4\sqrt{3}}{11}$, so our answer is $4+3+11=\boxed{018}$.

 



\textbf{Solution 6}

Let $CB = 1$, and let the feet of the altitudes from $D$ and $E$ to $\overline{CB}$ be $D'$ and $E'$, respectively. Also, let $DE = 8k$ and $EB = 15k$. We see that $BD' = 15k\cos B$ and $BE' = 23k\cos B$ by right triangles $\triangle{BDD'}$ and $\triangle{BEE'}$. From this we have that $D'E' = 8k\cos B$. With the same triangles we have $DD' = 23k\sin B$ and $EE' = 15k\sin B$. From 30-60-90 triangles $\triangle{CDD'}$ and $\triangle{CEE'}$, we see that $CD' = \frac{23k\sqrt{3}\sin B}{3}$ and $CE' = 15k\sqrt{3}\sin B$, so $D'E' = \frac{22k\sqrt{3}\sin B}{3}$. From our two values of $D'E'$ we get:$8k\cos B = \frac{22k\sqrt{3}\sin B}{3}$

 $\frac{\sin B}{\cos B} = \frac{8k}{\frac{22k\sqrt{3}}{3}} = \tan B$

 $\tan B = \frac{8}{\frac{22\sqrt{3}}{3}} = \frac{24}{22\sqrt{3}} = \frac{8\sqrt{3}}{22} = \frac{4\sqrt{3}}{11}$Our answer is then $4+3+11 = \boxed{018}$.

 


	\item (\textit{AIME I 2012}) Three concentric circles have radii $3,$ $4,$ and $5.$ An equilateral triangle with one vertex on each circle has side length $s.$ The largest possible area of the triangle can be written as $a + \frac{b}{c} \sqrt{d},$ where $a,$ $b,$ $c,$ and $d$ are positive integers, $b$ and $c$ are relatively prime, and $d$ is not divisible by the square of any prime. Find $a+b+c+d.$


 
 \textbf{Answer: }041



\textbf{Solution 1}

Reinterpret the problem in the following manner. Equilateral triangle $ABC$ has a point $X$ on the interior such that $AX = 5,$ $BX = 4,$ and $CX = 3.$ A $60^\circ$ counter-clockwise rotation about vertex $A$ maps $X$ to $X'$ and $B$ to $C.$ 

 
                    \begin{figure}[H]
                        \centering
                        \includegraphics[width=0.5\textwidth]{Solution/AIME_I_2012_13_Solution_Figure_1.png}
                    \end{figure}
                    

 Note that angle $XAX'$ is $60$ and $XA = X'A = 5$ which tells us that triangle $XAX'$ is equilateral and that $XX' = 5.$ We now notice that $XC = 3$ and $X'C = 4$ (recall that $C$ is the image of $B$ after rotation) which tells us that angle $XCX'$ is $90$ because there is a $3$-$4$-$5$ Pythagorean triple. Now note that $\angle ABC + \angle ACB = 120^\circ$ and $\angle XCA + \angle XBA = 90^\circ,$ so $\angle XCB+\angle XBC = 30^\circ$ and $\angle BXC = 150^\circ.$ Applying the law of cosines on triangle $BXC$ yields

 
 \[BC^2 = BX^2+CX^2 - 2 \cdot BX \cdot CX \cdot \cos150^\circ = 4^2+3^2-24 \cdot \frac{-\sqrt{3}}{2} = 25+12\sqrt{3}\]

 and thus the area of $ABC$ equals 
 \[\frac{\sqrt{3}}{4}\cdot BC^2 = 25\frac{\sqrt{3}}{4}+9.\]

 so our final answer is $3+4+25+9 = \boxed{041}.$

 



\textbf{Solution 2}

Here is a proof that shows that there are 2 distinct equilateral triangles (up to congruence) that have the given properties.

 We claim that there are 2 distinct equilateral triangles (up to congruence) that have the given properties; one of which has largest area. We have two cases to consider; either the center $O$ of the circles lies in the interior of triangle $ABC$ or it does not (and we shall show that both can happen). To see that the first case can occur proceed as follows. Using the notation from Solution 1, let $\angle XAC = \theta$ so that  $\angle BAX=60-\theta$. Let $AB=BC=AC=x$. The law of cosines on $\triangle BAX$ and $\triangle CAX$ yields 
 \[\cos(60-\theta)=\frac{x^2+9}{10x}\quad \text{ and }\quad  \cos\theta=\frac{x^2+16}{10x}.\] Solving this system will yield the value of $x$. Since $\cos\theta=\frac{x^2+16}{10x}$ we have that 
 \[\sin\theta=\frac{\sqrt{100x^2-(x^2+16)^2}}{10x}.\] Substituting these into the equation 
 \[\frac{x^2+9}{10x}=\cos(60-\theta)=\frac{1}{2}(\cos\theta+ \sqrt{3}\sin\theta)\] we obtain 
 \[\frac{x^2+9}{10x}=\frac{1}{2}\frac{x^2+16}{10x}+\frac{\sqrt{3}}{2}\frac{\sqrt{100x^2-(x^2+16)^2}}{10x}.\] After clearing denominators, combining like terms, isolating the square root, squaring, and expanding, we obtain 
 \[x^4-50x^2+193=0\] so that by the quadratic formula $x^2=25\pm12\sqrt{3}$. Under the hypothesis that $O$ lies in the interior of triangle $ABC$, $x^2$ must be $25+12\sqrt{3}$. To see this, note that the other value for $x^2$ is roughly $4.2$ so that $x\approx 2.05$, but since $AX=5$ and $AX\leq x$ we have a contradiction. We then obtain the area as in Solution 1.

 Now, suppose $O$ does not lie in the interior of triangle $ABC$. We then obtain convex quadrilateral $OBAC$ with diagonals $CB$ and $OA$ intersecting at $X$. Here $AX=AB=AC=x$. We may let $\alpha$ denote the measure of angle $CAX$ so that angle $XAB$ measures $60-\alpha$. Note that the law of cosines on triangles $CXA$ and $BXA$ yield the same equations as in the first case with $\theta$ replaced with $\alpha$. Thus we obtain again $x^2=25\pm12\sqrt{3}$. If $x^2=25+12\sqrt{3}$ then $x\approx 6.8$, but this is impossible since $AX\leq 5$ but the shortest possible distance from $A$ to $X$ is the height of equilateral triangle $ABC$ which is $\approx6.8\sqrt{3}\approx5.8$; a contradiction. Hence in this case $x^2=25-12\sqrt{3}$.  But, the area of this triangle is clearly less than that in the first case, so we are done. Hence the phrasing of the question (the triangle with maximal area) is absolutely necessary since there are 2 possible triangles (up to congruence).

 



\textbf{Solution 3}


                    \begin{figure}[H]
                        \centering
                        \includegraphics[width=0.5\textwidth]{Solution/AIME_I_2012_13_Solution_Figure_2.png}
                    \end{figure}
                    

 Let's call the circle center $X$. It has a distance of 3, 4, 5 to an equilateral triangle $LMN$. Consider $X$’s pedal triangle $ABC$. Since $X$’s antipedal triangle is equilateral, $X$ must be the one of the isogonic centers of $\triangle{ABC}$. We’ll take the one inside $ABC$, i.e., the Fermat point, because it leads to larger $\triangle LMN$. Now we construct the three equilateral triangles $ABD$, $ACE$, and $BCF$, the same way the Fermat point is constructed. Then we have $\angle DXE = \angle EXF = \angle FXE = 120$. Since $AEMCX$ is concyclic with $XM$=4 as diameter, we have $AC=4\sin(60)$. Similarly, $AB=3\sin(60)$, and $BC=5\sin(60)$. So $\triangle ABC$ is a 3-4-5 right triangle with $\angle BAC=90$. With some more angle chasing we get 
 \[\angle MXC+\angle LXB = \angle MAC + \angle LAB = 180 – \angle BAC = 90\]
 \[\angle LXM = 360 – (\angle MXC + \angle LXB + \angle BXC) = 360 –(90+120)=150\]By Law of Cosines, we have 
 \[LM^2 = 3^2+4^2-2*3*4\cos(150)=25+12\sqrt 3\]And the area follows.
 \[[LMN] = \frac{25}{4}\sqrt{3} + 9; \boxed{041}.\] By Mathdummy

 



\textbf{Solution 4}

Let $ABC$ be the equilateral triangle with $AB=BC=CA=x.$ Assume the coordinates of the vertices are $A(-\frac{x}{2},0), B(\frac{x}{2},0)$ and $C(0,\frac{\sqrt{3}}{2}x).$ Let $P(a,b)$ be such that $PA=3, PB=4$ and $PC=5.$ Then
 \[\left (a+\dfrac{x}{2}\right )^2 +b^2=3^2,\]
 \[\left (a-\dfrac{x}{2}\right )^2 +b^2=4^2,\]
 \[a^2+\left (b-\dfrac{\sqrt{3}}{2}x\right )^2 =5^2.\]Subtraction and addition of the first two equations yield $2ax=-7, 2a^2+\dfrac{1}{2}x^2+2b^2=25.$ The third equation gives $2a^2+2b^2-2\sqrt{3}bx+\dfrac{3}{2}x^2=50.$ Then $x^2-2\sqrt{3}bx=25.$ We can then solve for $a, b$ in terms of $x$ and have a substitution. We have 
 \[2\left (\dfrac{-7}{2x}\right )^2 + \dfrac{1}{2}x^2 + 2\left (\dfrac{x^2-25}{2\sqrt{3}x}\right )^2 =25.\]Simplify it we have a quadratic equation for $x^2: x^4-50x^2+193=0.$ So $x^2=25\pm 12\sqrt{3}.$ The larger one leads to the solution. The smaller one relates to another equilateral triangle, as indicated in Solution 2.

 -JZ

 



\textbf{Solution 5}


                    \begin{figure}[H]
                        \centering
                        \includegraphics[width=0.5\textwidth]{Solution/AIME_I_2012_13_Solution_Figure_3.png}
                    \end{figure}
                    We have $x=3$, $y=4$, and $z=5$. Because $AD=m$ is the median of $\triangle AOB$, by Stewart's Theorem we have 
 \[m^2=\frac 12 (x^2+y^2)-\frac 14 \cdot z^2\quad \Rightarrow \quad m = \frac 52.\] Because $CD$ is the altitude of equilateral triangle $OBC$, we have $CD=\frac{\sqrt{3}}2\cdot z$. Then in $\triangle ADC$, we have $\angle ADC=\varphi+90^\circ$, so $\cos(\angle ADC)=-\sin\varphi$, and the Law of Cosines gives
 \[s^2=m^2+\frac 34\cdot z^2 + mz\sqrt{3}\sin\varphi = 25 \left(1+\frac {\sqrt{3}}{2}\cdot\sin\varphi  \right)\]To calculate $\sin\phi$ we apply the Law of Cosines to $\triangle ADO$ to get
 \[mz\cos\varphi = m^2+\frac 14\cdot z^2 - x^2 \quad \Rightarrow \quad \cos\varphi = \frac 7{25}\quad \Rightarrow \quad \sin\varphi = \frac {24}{25}.\]Finally, we get $s^2=25+12\sqrt{3}$ and and thus the area of $\triangle ABC$ equals
 \[\frac{\sqrt{3}}{4}\cdot s^2 = 25\frac{\sqrt{3}}{4}+9.\]so our final answer is $3+4+25+9 = \boxed{041}.$ 

 \textbf{vladimir.shelomovskii@gmail.com, vvsss}

 



\textbf{Solution 6 (fastest)}

We note that this question is equivalent to finding the area of an equilateral triangle $ABC$ where there exists a point $P$ inside the triangle where $PA = 3$, $PB=4$, $PC=5$. Now, we reflect $P$ across $AB$, $BC$, and $CA$, call these points $D$, $E$, and $F$, respectively. It is obvious that hexagon $[ADBECF] = 2[ABC]$, so we try to find the area of this hexagon. Note that angle $DAF$ is equal to the sum of angles $DAP$ and $PAF$, which are equal to $2\angle{BAP}$ and $2\angle{PAC}$, respectively. Because $\angle{BAP} + \angle{PAC} = 60$, we have that $\angle{DAF} = 120$. By similar reasoning, we know $\angle{DAF} = \angle{EBD} = \angle{FCE} = 120$. Now, split this hexagon into 4 triangles, $DAF$, $EBD$, $FCE$, and $DEF$. It is easy to see that $DAF$, $EBD$, $FCE$ are 30-30-120 triangles with areas $\dfrac{9\sqrt{3}}{4}$, $4\sqrt{3}$, and $\dfrac{25\sqrt{3}}{4}$, respectively. We note that $DEF$ is a 3-4-5 right triangle, with side lengths $3\sqrt{3}$, $4\sqrt{3}$, and $5\sqrt{3}$, meaning that $DEF$ has area $18$.

 Summing, we have $[ADBECF] = 18 + \dfrac{25\sqrt{3}}{2}$

 Thus, $[ABC] = 9+\dfrac{25\sqrt{3}}{4}$, $9+25+3+4 = \boxed{41}$

 



\textbf{Solution 7 (faster than the fastest)}

Label the center of the circles $X$, and proceed with Solution 1 until out of the two possible equilateral triangles, you find the larger one.

 We note that the square of the side length $s$ of an equilateral triangle $ABC$ with the only information given being an arbitrary point $X$ inside $ABC$ such that $XA = a, XB = b, XB = c$ is given by the following formula:

 
 \[s^2 = \frac{1}{2}(a^2+b^2+c^2) + 2\sqrt{3}\delta,\]

 where $\delta$ equals the area of the triangle with side lengths $a,b,c.$ Since the $3,4,5$ triangle has area $6, \delta=6$ for us. Plugging in the other values $(XA = 5, XB = 4, XC = 3)$, we get

 
 \[s^2 = \frac{1}{2}(3^2+4^2+5^2) + 12\sqrt{3}.\]

 Multiplying this by $\frac{\sqrt{3}}{4},$ we attain the answer of 

 
 \[\frac{s^2\sqrt{3}}{4} = \frac{\sqrt{3}}{4}(25 + 12\sqrt{3}) = 9 + \frac{25\sqrt{3}}{4}, \Longrightarrow 9+25+3+4 = \boxed{041}.\]

 


	\item (\textit{AIME I 2012}) Complex numbers $a,$ $b,$ and $c$ are zeros of a polynomial $P(z) = z^3 + qz + r,$ and $|a|^2 + |b|^2 + |c|^2 = 250.$ The points corresponding to $a,$ $b,$ and $c$ in the complex plane are the vertices of a right triangle with hypotenuse $h.$ Find $h^2.$


 
 \textbf{Answer: }375



\textbf{Solution 1}

By Vieta's formula, the sum of the roots is equal to 0, or $a+b+c=0$. Therefore, $\frac{(a+b+c)}{3}=0$. Because the centroid of any triangle is the average of its vertices, the centroid of this triangle is the origin. Let one leg of the right triangle be $x$ and the other leg be $y$. Without the loss of generality, let $\overline{ac}$ be the hypotenuse. The magnitudes of $a$, $b$, and $c$ are just $\frac{2}{3}$ of the medians because the origin, or the centroid in this case, cuts the median in a ratio of $2:1$. So, $|a|^2=\frac{4}{9}\cdot((\frac{x}{2})^2+y^2)=\frac{x^2}{9}+\frac{4y^2}{9}$ because $|a|$ is two thirds of the median from $a$. Similarly, $|c|^2=\frac{4}{9}\cdot(x^2+(\frac{y}{2})^2)=\frac{4x^2}{9}+\frac{y^2}{9}$. The median from $b$ is just half the hypotenuse because the median of any right triangle is just half the hypotenuse. So, $|b|^2=\frac{4}{9}\cdot\frac{x^2+y^2}{4}=\frac{x^2}{9}+\frac{y^2}{9}$. Hence, $|a|^2+|b|^2+|c|^2=\frac{6x^2+6y^2}{9}=\frac{2x^2+2y^2}{3}=250$. Therefore, $h^2=x^2+y^2=\frac{3}{2}\cdot250=\boxed{375}$.

 



\textbf{Solution 2}

Note that by vieta's, the sum of the roots is $0,$ so the centroid (at (a+b+c)/3) of the right triangle is at the origin. This is all that is required by the structure of the polynomial. Let $A,B,C$ be the vertices of the triangle corresponding to roots $a,b,c$ with (WLOG) $\angle ACB=0.$ Let $O$ be the centroid of the triangle, and let $Z$ be the midpoint of $AB.$ By the property of the centroid as the intersection of the medians, $C,O,$ and $Z$ are collinear in that order. Also, by properties of medians and right triangles $AZ=BZ=CZ$ and $OC=2ZO.$ Then, set $AZ=BZ=CZ=3k$ and $OZ=k$ and $CO=2k.$ Stewart's theorem on ACZ and BCZ and rearranging gets
 \[3AO^2=12k^2+AC^2,\]
 \[3BO^2=12k^2+BC^2.\]Summing these up gives us 
 \[3(AO^2+BO^2)=24k^2+(AC^2+BC^2)\]Because $AC^2+BC^2=AB^2=36k^2,$ we have
 \[AO^2+BO^2=20k^2.\]Because $CO=2k,$
 \[AO^2+BO^2+CO^2=24k^2.\]By the second condition, this must be equal to $250.$ Then, $k^2=\frac{125}{12},$ and the hypotenuse squared is 
 \[(6k^2)^2=36k^2=36\left(\dfrac{125}{12}\right)=\boxed{375}.\] 

 



\textbf{Solution 3}

Assume $q$ and $r$ are real, so at least one of $a,$ $b,$ and $c$ must be real, with the remaining roots being pairs of complex conjugates. Without loss of generality, we assume $a$ is real and $b$ and $c$ are $x + yi$ and $x - yi$ respectively. By symmetry, the triangle described by $a,$ $b,$ and $c$ must be isosceles and is thus an isosceles right triangle with hypotenuse $\overline{bc}.$ Now since $P(z)$ has no $z^2$ term, we must have $a+b+c = a + (x + yi) + (x - yi) = 0$ and thus $a = -2x.$ Also, since the length of the altitude from the right angle of an isosceles triangle is half the length of the hypotenuse, $a-x=y$ and thus $y=-3x.$ We can then solve for $x$:

 
 \begin{align*} |a|^2 + |b|^2 + |c|^2 &= 250\\ |-2x|^2 + |x-3xi|^2 + |x+3xi|^2 &= 250\\ 4x^2 + (x^2 + 9x^2) + (x^2 + 9x^2) &= 250\\ x^2 &= \frac{250}{24} \end{align*}


 Now $h$ is the distance between $b$ and $c,$ so $h = 2y = -6x$ and thus $h^2 = 36x^2 = 36 \cdot \frac{250}{24} = \boxed{375.}$

 



\textbf{Solution 4 (Messy)}

Let the roots $a$, $b$, and $c$ each be represented by complex numbers $m + ni$, $p + qi$, and $r + ti$. By Vieta's formulas, their sum is 0. Breaking into real and imaginary components, we get:

 $m + p + r = 0$

 $n + q + t = 0$

 And, we know that the sum of the squares of the magnitudes of each is 250, so

 $m^2 + n^2 + p^2 + q^2 + r^2 + t^2 = 250$

 Given the complex plane, we set each of these complex numbers to points: $(m, n)$, $(p, q)$, $(r, t)$. WLOG let $(r, t)$ be the vertex opposite the hypotenuse.

 If the three points form a right triangle, the vectors from $(r, t)$ to $(m, m)$ and $(p, q)$'s dot product is 0.$mp + r^2 - r(m + p) + nq + t^2 - t(n + q) = 0$

 Substituting $m + p + r = 0$ and likewise, simplifying: $mp + 2r^2 + nq + 2t^2 = 0$

 Rearranging we get:$r^2 + t^2 = -\frac{mp + nq}{2}$

 The answer is the distance from $(m, n)$ to $(p, q)$ = $m^2 + n^2 + p^2 + q^2 - 2(mp + nq)$.Substituting the equation equal to 250,

 $= 250 - r^2 - t^2 - 2(mp + nq)$$= 250 + \frac{mp + nq}{2} - 2(mp + nq)$$= 250 - \frac{3}{2} \cdot (mp + nq)$

 Taking our original equations summing to 0, and squaring each we get:

 $n + q = -t$; $m + p = -r$

 $n^2 + 2nq + q^2 = t^2$; $m^2 + 2mp + p^2 = r^2$

 Adding, we get: 

 $m^2 + n^2 + p^2 + q^2 + 2(mp + nq) = r^2 + t^2$

 Substituting again we obtain:

 $250 - r^2 - t^2 + 2(mp + nq) = r^2 + t^2$$2(r^2 + t^2) = 250 + 2(mp + nq)$$r^2 + t^2 = 125 + (mp + nq)$

 Substituting the equivalence of $r^2 + t^2$:

 $-\frac{mp + nq}{2} = 125 + (mp + nq)$

 Solving for $mp + nq$, we find it equal to $-\frac{250}{3}$.

 Substituting this value into our answer expression, we get:

 $250 - \frac{3}{2} \cdot (-\frac{250}{3})$, Answer = $\boxed{375}$.

 



\textbf{Solution 5 (clean)}

As noted in the previous solutions, $a+b+c = 0$.Let $a = a_1+a_2 i$, $b = b_1+b_2 i$, $c = c_1+c_2 i$ and we have $a_1 + b_1 + c_1 = a_2 + b_2 + c_2 = 0$. Then the given $|a|^2 + |b|^2 + |c|^2 = 250$ translates to $\sum_{} ( {a_1}^2 + {a_2}^2 ) = 250.$Note that in a right triangle, the sum of the squares of the three sides is equal to two times the square of the hypotenuse, by the pythagorean theorem.Thus, we have 
 \[2h^2 = (a_1 - b_1)^2 + (a_2 - b_2)^2 + (b_1 - c_1)^2 + (b_2 - c_2)^2 + (a_1 - c_1)^2 + (a_2 - c_2)^2\] 
 \[= 2 \left( \sum_{} ( {a_1}^2 + {a_2}^2 )  \right) - 2 \left( \sum_{cyc} a_1 b_1 + \sum_{cyc} a_2 b_2 \right)\]
 \[= 500 - \left( (a_1 + b_1 + c_1)^2 + (a_2 + b_2 + c_2)^2 -  \sum_{cyc} ( {a_1}^2 + {a_2}^2 )     \right)\]
 \[= 500 - (0^2 + 0^2 - 250)\]so $h^2 = \boxed{375}$ and we may conclude.~ rzlng

 



\textbf{Solution 5 (dirty because of god intuition + wishful thinking)}

First, note that the roots of this cubic will be $a, b$ and $-(a+b)$ due to Vieta's, which means that the sum of the roots are 0.

 Now, we use some god level wishful thinking. It would really be nice if one of these roots was on the real axis, and it was a right isosceles triangle, because that would be very convenient and easy to work with. The neat part is that it actually works

 Set one of the roots as $r$, where $r$ is any real number. WLOG, assume that this is the right angle. With symmetry to respect of the x axis (because symmetry makes the imaginary parts of the other 2 roots cancel out, besides the fact that complex conjugate root theorem forces it). This way, we can set the other 2 roots as $\frac{r}{2}+ni$ and $\frac{r}{2}-ni$, making the roots add up to 0. 

 Now, as we want the roots to satisfy the original condition (right triangle) we are going to have to set an equation to find $n$ out. We use the fact that it is an isosceles right triangle to find that $\frac{3r}{2}=n$, which means that the 2 other roots are now $\frac{r}{2}+\frac{3r}{2}i$ and $\frac{r}{2}-\frac{3r}{2}i$

 Now we use the fact that $|a|^2+|b|^2+|c|^2=250$. Clearly one of these is $r$ away from the origin, so that gets $r^2$, and then we get $2*\frac{r}{2}^2+\frac{3r}{2}^2$ which gets us $5r^2$, getting $r^2+5r^2=250$, so $r=\sqrt{\frac{250}{6}}$. So the final answer comes out to
 \[(\frac{250}{6}*\frac{9}{4}*2)^2=\boxed{\boxed{375}}\]

 -dragoon

 P.S.: The main propose, saying that one root is on the real axis and the right triangle is a right isosceles triangle is not actually only a wishful thinking that came out by luck, but actually is something that must be true due to the complex conjugate root theorem. -SuperDolphin

 - in this specific problem this logic is incorrect, as the coefficients of the polynomial can be complex. If a question specified real roots, then this would be appropriate. In this case, it is just wishful thinking as stated above

 



\textbf{Solution 6 (vectors)}

As shown in the other solutions, $a+b+c = 0$.

 Without loss of generality, let $b$ be the complex number opposite the hypotenuse.

 Note that there is an isomorphism between $\mathbb{C}$ under $+$ and $\mathbb{R}^2$ under $+$.

 Let $\Vec{a}$, $\Vec{b}$, and $\Vec{c}$ be the corresponding vectors to $a$, $b$, and $c$.

 Thus $\Vec{a} + \Vec{b} + \Vec{c} = \Vec{0}$

 $\Rightarrow 0 = \Vec{0}\cdot \Vec{0} = (\Vec{a} + \Vec{b} + \Vec{c})\cdot (\Vec{a} + \Vec{b} + \Vec{c}) = \Vec{a}\cdot \Vec{a} + \Vec{b}\cdot \Vec{b} + \Vec{c}\cdot \Vec{c} + 2(\Vec{a}\cdot \Vec{b} + \Vec{a}\cdot \Vec{c} + \Vec{b}\cdot \Vec{c})$

 Now $|a|^2 + |b|^2 + |c|^2 = 250$ implies that $\lVert \Vec{a}\rVert^2 + \lVert \Vec{b}\rVert^2 + \lVert \Vec{c}\rVert^2 = 250$

 $\Rightarrow \Vec{a}\cdot \Vec{a} + \Vec{b}\cdot \Vec{b} + \Vec{c}\cdot \Vec{c} = \lVert \Vec{a}\rVert^2 + \lVert \Vec{b}\rVert^2 + \lVert \Vec{c}\rVert^2 = 250$

 Also note that because there is a right angle at $b$, $\Vec{a} - \Vec{b}$ and $\Vec{c} - \Vec{b}$ are perpendicular.

 $\Rightarrow (\Vec{a} - \Vec{b})\cdot (\Vec{c} - \Vec{b}) = 0$

 $\Rightarrow 0 = (\Vec{a} - \Vec{b})\cdot (\Vec{c} - \Vec{b}) = \Vec{a}\cdot \Vec{c} + \Vec{b} \cdot \Vec{b} - \Vec{a} \cdot \Vec{b} - \Vec{b} \cdot \Vec{c}$

 Note that $h^2 = |a-c|^2$

 $\Rightarrow h^2 = \lVert \Vec{a} - \Vec{c} \rVert^2 = (\Vec{a} - \Vec{c})\cdot (\Vec{a} - \Vec{c}) = \Vec{a} \cdot \Vec{a} + \Vec{c}\cdot \Vec{c} - \Vec{a}\cdot \Vec{c} - \Vec{a}\cdot \Vec{c} = \Vec{a} \cdot \Vec{a} + \Vec{c}\cdot \Vec{c} - 2\Vec{a}\cdot \Vec{c}$.

 $\Rightarrow 0 = \Vec{a}\cdot \Vec{a} + \Vec{b}\cdot \Vec{b} + \Vec{c}\cdot \Vec{c} + 2(\Vec{a}\cdot \Vec{b} + \Vec{a}\cdot \Vec{c} + \Vec{b}\cdot \Vec{c}) = 250 + 2(\Vec{a}\cdot \Vec{b} + \Vec{a}\cdot \Vec{c} + \Vec{b}\cdot \Vec{c})$

 $\Rightarrow -250 = 2(\Vec{a}\cdot \Vec{b} + \Vec{a}\cdot \Vec{c} + \Vec{b}\cdot \Vec{c})$

 $\Rightarrow -125 = \Vec{a}\cdot \Vec{b} + \Vec{a}\cdot \Vec{c} + \Vec{b}\cdot \Vec{c}$

 $\Rightarrow -125 = -125 + 0 = (\Vec{a}\cdot \Vec{b} + \Vec{a}\cdot \Vec{c} + \Vec{b}\cdot \Vec{c}) + (\Vec{a}\cdot \Vec{c} + \Vec{b} \cdot \Vec{b} - \Vec{a} \cdot \Vec{b} - \Vec{b} \cdot \Vec{c}) = 2\Vec{a}\cdot \Vec{c} + \Vec{b} \cdot \Vec{b}$

 $\Rightarrow 125 = - \Vec{b} \cdot \Vec{b} - 2\Vec{a}\cdot \Vec{c}$

 $\Rightarrow 375 = 250 + 125 = \Vec{a}\cdot \Vec{a} + \Vec{b}\cdot \Vec{b} + \Vec{c}\cdot \Vec{c} - \Vec{b} \cdot \Vec{b} - 2\Vec{a}\cdot \Vec{c} = \Vec{a}\cdot \Vec{a} + \Vec{c}\cdot \Vec{c} - 2\Vec{a}\cdot \Vec{c} = h^2$

 $\Rightarrow h^2 = \boxed{375}$

 



\textbf{Solution 6 (similar to solution 4)}

Note that the roots of the polynomial $(a,b,c)$ must sum to $0$ due to the $z^2$ coefficient equaling $0$ because of Vieta's Formulas. This tells us that $a+b+c = 0$ and $\overline{a} + \overline{b} + \overline{c} = 0$ so $(a+b+c)(\overline{a+b+c}) = |a|^2 + |b|^2 + |c|^2 + a\overline{b} + \overline{a}b + b\overline{c} + \overline{b}c + c\overline{a} + \overline{c}a = 250 + a\overline{b} + \overline{a}b + b\overline{c} + \overline{b}c + c\overline{a} + \overline{c}a =  0$ so $a\overline{b} + \overline{a}b + b\overline{c} + \overline{b}c + c\overline{a} + \overline{c}a = -250.$

 Note that terms similar to $\overline{a}b + \overline{b}a$ appear in $|a-b|^2$ so we decide to sum $|a-b|^2 + |b-c|^2 + |c-a|^2$ out of intuition. Note that this corresponds to the sum of the squares of the sidelengths of the right triangle and if WLOG the $|c-a|^2$ side is the hypotenuse then our sum is equal to $2|c-a|^2 = 2h^2.$

 
 \[2|a|^2 + 2|b|^2 + 2|c|^2 -a\overline{b} - \overline{a}b - b\overline{c} - \overline{b}c - c\overline{a} - \overline{c}a = 500 + 250 = 2h^2\]

 As a result we know that $h^2 = \boxed{375}$~SailS

 



\textbf{Solution 7}

By Vieta's formula, $a+b+c=0$. Let $a=x_1+y_1i,b=x_2+y_2i,c=x_3+y_3i$. Adding and equating the real part and the imaginary part with $0$, we have\begin{gather*}x_1+x_2+x_3=0 \tag{1}\\y_1+y_2+y_3=0 \tag{2}\end{gather*}It is given that $x_1^2+x_2^2+x_3^2+y_1^2+y_2^2+y_2^2=250$. By symmetry, we can construct the right triangle $\triangle ABC$, which satisfies $\angle B=90^\circ$, with the following coordinates:\begin{align*}A&=(x_1,y_1)\\B&=(x_2,y_2)\\C&=(x_3,y_3)\end{align*}
We have established $\overline{AB}\perp \overline{BC}$, which means the dot product of vectors $\vec{u}=(x_2-x_1, y_2-y_1)$ and $\vec{v}=(x_3-x_2, y_3-y_2)$ evaluates to $0$:\begin{gather*}(x_2-x_1)(x_3-x_2)+(y_2-y_1)(y_3-y_2)=0\\x_1x_2+x_2x_3-x_2^2+x_1x_3+y_1y_2+y_2y_3-y_2^2+y_1y_3=0\\\end{gather*}Simplifying with $x_1+x_3=-x_2$ and $y_1+y_3=-y_2$, the above equation is equivalent to\begin{gather*}2(x^2+y^2)=-(x_1x_3+y_1y_3) \tag{3}\end{gather*}Squaring $(1)$ and $(2)$ and adding, we have
 \[x_1^2+x_2^2+x_3^2+y_1^2+y_2^2+y_2^2+2(x_1x_2+x_2x_3+x_3x_1+y_1y_2+y_2y_3+y_3y_1)=0\]Or
 \[250-2(x_2^2+y_2^2)+2(x_3x_1+y_3y_1)=0\]Substituting $(3)$ gives $x_2^2+y_2^2=\frac{125}{3}$.

 Solving for $AC^2$, we see that $h^2=(x_3-x_1)^2+(y_3-y_1)^2=250+3(x_2^2+y_2^2)=\boxed{375}$.

 


	\item (\textit{AIME I 2012}) There are $n$ mathematicians seated around a circular table with $n$ seats numbered $1,$ $2,$ $3,$ $...,$ $n$ in clockwise order. After a break they again sit around the table. The mathematicians note that there is a positive integer $a$ such that


 <ul>
($1$) for each $k,$ the mathematician who was seated in seat $k$ before the break is seated in seat $ka$ after the break (where seat $i + n$ is seat $i$);
</ul><ul>
($2$) for every pair of mathematicians, the number of mathematicians sitting between them after the break, counting in both the clockwise and the counterclockwise directions, is different from either of the number of mathematicians sitting between them before the break.
</ul>Find the number of possible values of $n$ with $1 < n < 1000.$


 
 \textbf{Answer: }332



\textbf{Solution}

It is a well-known fact that the set $0, a, 2a, ... (n-1)a$ forms a complete set of residues if and only if $a$ is relatively prime to $n$.

 Thus, we have $a$ is relatively prime to $n$. In addition, for any seats $p$ and $q$, we must have $ap - aq$ not be equivalent to either $p - q$ or $q - p$ modulo $n$ to satisfy our conditions. These simplify to $(a-1)p \not\equiv (a-1)q$ and $(a+1)p \not\equiv (a+1)q$ modulo $n$, so multiplication by both $a-1$ and $a+1$ must form a complete set of residues mod $n$ as well. 

 Thus, we have $a-1$, $a$, and $a+1$ are relatively prime to $n$. We must find all $n$ for which such an $a$ exists. $n$ obviously cannot be a multiple of $2$ or $3$, but for any other $n$, we can set $a = n-2$, and then $a-1 = n-3$ and $a+1 = n-1$. All three of these will be relatively prime to $n$, since two numbers $x$ and $y$ are relatively prime if and only if $x-y$ is relatively prime to $x$. In this case, $1$, $2$, and $3$ are all relatively prime to $n$, so $a = n-2$ works.

 Now we simply count all $n$ that are not multiples of $2$ or $3$, which is easy using inclusion-exclusion. We get a final answer of $998 - (499 + 333 - 166) = \boxed{332}$.

 
Note: another way to find that $(a-1)$ and $(a+1)$ have to be relative prime to $n$ is the following: start with $ap-aq \not \equiv \pm(p-q) \pmod n$. Then, we can divide by $p-q$ to get $a \not \equiv \pm 1$ modulo $\frac{n}{\gcd(n, p-q)}$. Since $\gcd(n, p-q)$ ranges through all the divisors of $n$, we get that $a \not \equiv \pm 1$ modulo the divisors of $n$ or $\gcd(a-1, n) = \gcd(a+1, n) = 1$.

 



\textbf{Solution 2}

We see that $\gcd(a, n)=1$. Let $f(x)=ax\bmod{n}$ detote the seat number after the break.The second condition implies that for $b,c\in\{1,2,\ldots, n\}, b\neq c$, we have\begin{align*}\lvert b-c\rvert -1 &\neq \lvert f(b)-f(c)\rvert -1\\\lvert b-c\rvert -1 &\neq n- \lvert f(b)-f(c)\rvert -1\\n-\lvert b-c\rvert -1 &\neq n- \lvert f(b)-f(c)\rvert -1\\n-\lvert b-c\rvert -1 &\neq n- \lvert f(b)-f(c)\rvert -1\end{align*}
Or\begin{gather}\lvert b-c\rvert \neq \lvert a(b-c)\bmod{n}\rvert\\\lvert b-c\rvert \neq n - \lvert a(b-c)\bmod{n}\rvert\\\end{gather}WLOG let $u = (b-c)\in\{1,2,\ldots n-1\}$.

 \textbf{Case 1: }$\gcd{(u,n)}=1$

 The integers $a$ that violates the equations above satisfy the following:
 \[a\equiv 1, -1\bmod{n}\]

 \textbf{Case 2: }$\gcd{(u,n)}=d>1$

 Let $u=d\cdot t_1, n=d\cdot t_2, \gcd(t_1,t_2)=1$. The undesired $a$ satisfy\begin{align*}d\cdot t_1=ad\cdot t_1\bmod{d\cdot t_2}\\d\cdot t_1=d\, (at_1\bmod{t_2})\\0=(a-1)t_1\bmod{t_2}\end{align*}
This means $a\equiv 1\bmod{t_2}$. Similarly, we obtain $a\equiv -1\bmod{t_2}$ from $(2)$.

 Let $d_i\ge 2$ be a positive divisor of $n$. We want the number of valid $n$ such that there exists an integer $a$ satisfying
 \[a\not\equiv -1, 0, 1 \pmod{d_i}\]Clearly no such $a$ exists when $2\mid n$ or $3\mid n$. Otherwise, the system\begin{align*}a&\equiv 2,3,\ldots d_1 -2\bmod{d_1}\\a&\equiv 2,3,\ldots d_2 -2\bmod{d_2}\\&\ldots\\a&\equiv 2,3,\ldots d_k -2\bmod{d_k}\\\end{align*}
guarantees a solution for $a$ modulo $n$ residue that is relatively prime with $n$, such as $a=2$. 

 Counting the integers $n\in\{2,3,\ldots 999\}$ that are divisible by neither $2$ nor $3$, we have
 \[998-499-333+166=\boxed{332}\]

 


\end{enumerate}